\documentclass[11pt]{article}
\usepackage[margin=1in]{geometry}
\usepackage{amsmath,amssymb,amsthm}
\title{Disproof of Conjecture 00000005344}
\author{earthking11}
\date{October 6, 2026}
\newtheorem{proposition}{Proposition}
\begin{document}
\maketitle
\section*{Counterexample}
Every ordinary category can be regarded as a strict bicategory whose only
$2$-morphisms are identities.  In this locally discrete bicategory,
horizontal composition of $1$-morphisms is exactly ordinary composition.
Thus a category with noncommuting endomorphisms is enough to refute strict
commutativity.

Use the category of sets and the Boolean set $B=\{\mathrm{false},
\mathrm{true}\}$.  Define endomorphisms
\[
 f(x)=\mathrm{false},\qquad g(x)=\neg x.
\]
At $x=\mathrm{false}$ the two composite orders give
\[
 (f\circ g)(x)=\mathrm{false},\qquad
 (g\circ f)(x)=\mathrm{true}.
\]
Hence $f\circ g\ne g\circ f$.

\begin{proposition}
Horizontal composition in a bicategory is not strictly commutative in
general.
\end{proposition}
\begin{proof}
Regard the category of sets as the locally discrete strict bicategory just
described.  The Boolean endomorphisms $f$ and $g$ have unequal horizontal
composites, as witnessed at $\mathrm{false}$.
\end{proof}

The bicategory axioms provide associativity and units (up to coherent
$2$-isomorphism in the general case); they do not provide commutativity.
Nor can normalizing a numerical constant to one turn two unequal
$1$-morphisms into equal ones.

The Lean project kernel-checks the two function evaluations, their inequality,
and the negation of universal strict commutativity.
\end{document}
